% ======================================================================
%  poster-research-terraces.tex -- START HERE. Everything visual lives in postertheme.sty;
%  this file is just your content. Fill in \PosterSetup, then write
%  (or paste) your sections, theorems, and figures below.
%
%  COMPILE WITH LuaLaTeX:
%    * Overleaf:  Menu -> Compiler -> LuaLaTeX, then Recompile.
%    * Locally:   latexmk
% ======================================================================
\DocumentMetadata{lang=en-US, pdfstandard=ua-1, tagging=on}
\documentclass[12pt]{article}
\usepackage{postertheme}

\PosterSetup{
  title      = {Calabi--Yau Manifolds},
  subtitle   = {What they are, why they exist, and why they matter},
  authors    = {Franklin Wang},
  department = {Department of Mathematics and Statistics, Auburn University},
  % the soul of the poster -- one sentence, read before everything else:
  highlight  = {One cheap topological test --- $c_1 = 0$ --- buys the richest geometry of all: $\operatorname{Ric} = 0$, guaranteed.},
  contact    = {zzw0109@auburn.edu},
  repo-url   = {https://github.com/Franklinwang72/auburn-university-latex-template},
  repo-name  = {github.com/Franklinwang72},
  qr-url     = {https://www.linkedin.com/in/franklin-wang-a14b7b327},
  qr-label   = {Connect on LinkedIn},
  % PDF metadata: type REAL dashes here (- -- do not turn into dashes
  % inside PDF strings)
  pdf-title  = {Calabi–Yau Manifolds — What They Are, Why They Exist, and Why They Matter},
  pdf-author = {Franklin Wang},
}

\begin{document}
\PosterHeader

\begin{postercolumns}
% ================= LEFT COLUMN =================
\begin{postercolumn}
% ---------- LEFT: the prerequisites, built up for an undergraduate ----------
\section{What Is a Manifold?}
A \auterm{manifold} is a mathematician's word for a smooth shape --- of
any dimension. Its one defining trick: up close it looks \auterm{flat},
the way the ground looks flat from where you stand even though the
Earth is round. The number of coordinates you need to pin down a point
is the \auterm{dimension}, and the shapes here are \auterm{compact}:
finite and edgeless, closed up like a sphere rather than sprawling like
an open plain.
\vfill
\section{Complex Dimensions}
Calabi--Yau manifolds run on \auterm{complex} coordinates. A complex
number $x + iy$ is a single arrow in the plane --- length and
direction bundled into one --- so one complex dimension already
carries two real ones. An $n$-complex-dimensional shape is therefore
$2n$-real-dimensional: a Calabi--Yau \auterm{threefold} $(n = 3)$ is a
\auterm{six}-dimensional world.
\vfill
\section{Curvature: Do Triangles Behave?}
How can you tell a surface is curved without stepping outside it? Draw
a triangle. On a \auterm{sphere} its angles sum to \auterm{more} than
$180^\circ$; on a \auterm{saddle}, to \auterm{less}; only on a
\auterm{flat} sheet do they hit $180^\circ$ exactly. An ant that never
leaves the surface could discover its world is round this way ---
curvature is \auterm{intrinsic}, and it can change from point to point.
\begin{center}
\includegraphics[width=0.94\textwidth,
  alt={Three small surfaces in a row: a blue sphere with a triangle of
    three outward-bulging arcs, a flat gray sheet with a
    straight-edged triangle, and a burnt-orange saddle surface with a
    triangle whose sides pinch inward.}]{figures/curvature-trio.png}
\par\vspace{0.15em}
{\PosterCaption\color{AUgray}
Sphere: more than $180^\circ$ \textperiodcentered\;
flat sheet: exactly $180^\circ$ \textperiodcentered\;
saddle: less than $180^\circ$.\par}
\end{center}
\vfill
\section{Ricci Curvature}
Past two dimensions, curvature splits into many independent pieces.
The \auterm{Ricci curvature} is the piece that controls \auterm{volume}.
Sprinkle a tiny ball of dust and let each speck drift along a straight
line: on a sphere the ball shrinks, on a saddle it swells, and when
$\operatorname{Ric} = 0$ it holds its volume exactly. Such a space is
genuinely curved yet volume-balanced --- which is precisely Einstein's
equation for a spacetime with no matter. \auterm{Ricci-flat spaces are
gravity's vacuums.}
\end{postercolumn}
% ================ CENTRE COLUMN ================
\begin{postercolumn}[wide]
% ---------- CENTER: the definition, the theorem, the picture ----------
\section{What Is a Calabi--Yau Manifold?}
A \auterm{Calabi--Yau manifold} is a compact complex manifold with two
features. \auterm{(1) K\"ahler:} it carries both a way to measure
distance --- a metric --- and complex coordinates, fitting together as
smoothly as possible. \auterm{(2) A holomorphic volume form} $\Omega$:
a way to gauge complex volume, built from the coordinates, that
vanishes \emph{nowhere} (in bundle language, a trivial canonical
bundle $K_X \cong \mathcal{O}_X$). Feature (2) forces a purely
\auterm{topological} quantity to vanish --- the \auterm{first Chern
class} $c_1(X)$, a single whole-number invariant of the shape you can
compute without measuring anything. Calabi asked whether that cheap
test forces a Ricci-flat metric to exist.
\begin{theorem}[Yau 1978 --- the Calabi Conjecture]
Let $(X,\omega)$ be a compact K\"ahler manifold with
$c_1(X) = 0$ in $H^2(X,\mathbb{R})$. Then $X$ carries a
\textit{unique} Ricci-flat K\"ahler metric $\tilde\omega$ with
$[\tilde\omega] = [\omega]$:
\punchline{\operatorname{Ric}(\tilde\omega) \;=\; 0 .}
Every Calabi--Yau manifold admits Ricci-flat metrics ---
exactly one in each K\"ahler class.
\end{theorem}
\vfill
\auterm{How is it proved?}\ Fix an overall scale; every candidate
metric is then $\omega + i\,\partial\bar\partial\varphi$ for a single
function $\varphi$, and Ricci-flatness becomes one nonlinear PDE ---
the \auterm{complex Monge--Amp\`ere equation}
% size change OUTSIDE the environment: \fontsize inside math is a no-op
{\fontsize{28}{33}\selectfont
\begin{equation*}
  \bigl(\omega + i\,\partial\bar\partial\varphi\bigr)^{\,n}
  \;=\; e^{F}\,\omega^{n}
\end{equation*}}
--- whose unknown sits inside a determinant of its own second
derivatives. Yau's a~priori estimates force a solution to exist.
\vfill
\begin{center}
\PosterFigure{A three-quarter view of a complex-curve slice of the
  Fermat quintic Calabi-Yau threefold: interlocking petal-like sheets
  that flow from Auburn orange at the top through pale gray to deep
  navy at the base, with a fivefold pinwheel symmetry and a glossy
  finish.}{%
  \includegraphics[width=0.43\textwidth]{figures/calabi-yau-quintic.png}}
\par\vspace{0.25em}
% caption constrained to the figure width so the two read as one block
\begin{minipage}{0.60\textwidth}
{\PosterCaption\color{AUgray}
A Calabi--Yau threefold has six real dimensions --- too many to draw,
but not too many to slice. Shown: the surface $z_1^5+z_2^5=1$, a
complex curve inside the Fermat quintic
$\,z_0^5+\dots+z_4^5=0\,$ in $\mathbb{P}^4$, projected to $3$-space;
its $25$ patches carry the standard Hanson parametrization.\par}
\end{minipage}
\end{center}
\end{postercolumn}
% ================= RIGHT COLUMN ================
\begin{postercolumn}
% ---------- RIGHT: examples a lay reader can hold, then the payoff ----------
\section{The Simplest One: A Donut}
In one complex dimension the definition selects exactly one shape: the
torus $\mathbb{C}/\Lambda$, the surface of a donut. Its volume form is
$dz$, inherited from the plane, and its Ricci-flat metric is the flat
geometry a sheet of paper keeps when rolled up in two directions ---
everything Yau's theorem promises, visible by hand.
\begin{center}
\includegraphics[width=0.27\textwidth,
  alt={A glossy slate-blue torus, the surface of a donut, with one
    vertical circle around the tube and one horizontal circle around
    the hole drawn in deep navy: the two directions in which the flat
    sheet was rolled up.}]{figures/cy-torus.png}
\par\vspace{0.1em}
{\PosterCaption\color{AUgray}
The torus: a flat sheet rolled up in two directions.\par}
\end{center}
\vfill
\section{The Family Album}
One dimension up sit the celebrated \auterm{K3 surfaces}; in dimension
three, the \auterm{quintic}. A shape cut out by one degree-$d$
equation in complex projective space $\mathbb{P}^{n+1}$ is
Calabi--Yau exactly when $d = n+2$:
\begin{center}
\PosterSmall
% tag the first row as TH cells (screen readers announce column headers)
\tagpdfsetup{table/header-rows={1}}%
\begin{tabular}{c l l}
  \toprule
  \auhead{dim} & \auhead{Calabi--Yau} & \auhead{why $K_X \cong \mathcal{O}_X$} \\
  \midrule
  $1$ & elliptic curve $\mathbb{C}/\Lambda$ & $dz$ descends \\
  $2$ & K3 surface, e.g.\ $X_4 \subset \mathbb{P}^3$ & $4 = 2+2$ \\
  $3$ & quintic $X_5 \subset \mathbb{P}^4$ & $5 = 3+2$ \\
  \bottomrule
\end{tabular}
\end{center}
\vspace{0.2em}
The quintic alone moves in a $101$-parameter family of Ricci-flat
worlds.
\vfill
\section{Six Hidden Dimensions}
String theory needs spacetime to have \auterm{ten} dimensions --- four
we see, six curled up too small to notice. In 1985 Candelas, Horowitz,
Strominger and Witten showed that when the hidden six form a
Calabi--Yau threefold, the physics left in four dimensions keeps
supersymmetry, with Yau's Ricci-flat metric as its vacuum. If string
theory is right, a shape like the one at the centre of this poster
sits at every point of space.
\vfill
\section{A Shape and Its Mirror}
Calabi--Yau threefolds come in pairs: each $X$ has a \auterm{mirror}
$X^{\vee}$ with its two ``shape numbers'' swapped --- one counts ways
to resize the shape, the other ways to deform it. String theory cannot
tell the pair apart, and in 1991 Candelas--de~la~Ossa--Green--Parkes
used this to count rational curves on the quintic ($2875$ lines,
$609{,}250$ conics) by an integral on the mirror.
\vfill
\accentsection{Takeaway}
{\PosterSemibold\fontsize{23.8}{32.5}\selectfont\raggedright
 \hyphenpenalty=10000\exhyphenpenalty=10000
A Calabi--Yau manifold passes one topological test, $c_1 = 0$ ---
and Yau's theorem rewards it with perfect geometry,
$\operatorname{Ric} = 0$: the meeting point of analysis, algebra, and
physics.\par}
\end{postercolumn}
\end{postercolumns}

\begin{posterfooter}
  \textbf{References:}\;
  E.~Calabi, Proc.\ ICM 1954 \textperiodcentered\;
  S.-T.~Yau, Comm.\ Pure Appl.\ Math.\ 31 (1978) \textperiodcentered\;
  P.~Candelas, X.~de~la~Ossa, P.~Green, L.~Parkes, Nucl.\ Phys.\ B 359 (1991) \textperiodcentered\;
  D.~Huybrechts, \emph{Complex Geometry}, Springer, 2005.\quad
  \textbf{Contact:}\; \href{mailto:\PosterContact}{\PosterContact}\quad
  \textbf{Template:}\; \href{\PosterURL}{\PosterURLShort}
\end{posterfooter}
\end{document}
